\section{The shape strip and its intrinsic coordinate}
\label{sec:shapes}

Put
\[
 H=h_X(0)=h_L(0),\qquad D_H=\{s\in\C:|\Im s|<H\},
\]
where \(H=\infty\) is allowed. The meromorphic curve on this strip is
\(\Gamma(s)=X(s,0)\). By \Cref{cor:shape-natural-boundary}, if
\(H<\infty\), the primary curve has no meromorphic continuation through
any point of either boundary line of \(D_H\). We shall first dispose of
its isotropic tangencies. In their absence we shall construct a global
shape coordinate for every continued impact, including impacts at poles
of the affine curve.

\subsection{The meromorphic normal and isotropic centers}

Use isotropic affine coordinates
\[
 Z=x_1+\ii x_2,\qquad W=x_1-\ii x_2.
\]
On the original real curve put
\[
 a=e^{\ii\phi},\qquad f=\phi',\qquad
 c=r(\phi)f^3>0.
\]
The normalization of the phase coordinate gives
\begin{equation}
 \Gamma'=\frac{c}{f^2}e_\phi,\qquad
 Z'=\frac{\ii c a}{f^2},\qquad
 W'=-\frac{\ii c}{a f^2}.
 \label{shape:coordinate-derivatives}
\end{equation}
In particular,
\begin{equation}
 a^2=-\frac{Z'}{W'},\qquad
 f=\frac{1}{2\ii}\frac{(a^2)'}{a^2},\qquad
 a=\frac{f^2Z'}{\ii c}.
 \label{shape:normal-reconstruction}
\end{equation}
Read the first expression as a meromorphic function, then use its
logarithmic derivative in the second expression, and finally use the
third expression to choose the normal itself. These formulas extend
\(a\) and \(f\) meromorphically throughout \(D_H\); all their identities
with the original real data continue there. We call a zero or a pole of
\(a\) an \emph{isotropic center}.

We shall repeatedly use the intrinsic normalization
\begin{equation}
 \frac{\det(\Gamma',\Gamma'')^4}
 {(\Gamma'\mathbin{\cdot}\Gamma')^3}=c^2.
 \label{shape:intrinsic}
\end{equation}
This is the identity from \Cref{cont:intrinsic-normalization}. It also
follows immediately from \eqref{shape:coordinate-derivatives}:
\(\Gamma'\cdot\Gamma'=c^2/f^4\) and
\(\det(\Gamma',\Gamma'')=c^2/f^3\). Under a reparametrization
\(s=\sigma(u)\), the left side of \eqref{shape:intrinsic} is multiplied
by \(\sigma'(u)^6\).

\begin{lemma}[Rational profiles at an isotropic center]
\label{shape:isotropic-profile}
Suppose that \(a\) has a zero of order \(n\geq1\) at \(p\in D_H\).
Then \(n\ne3\). Set
\[
 D_+=n+3,\qquad D_-=3-n.
\]
There is a germ \(\xi(w)=w+O(1)\) at infinity such that both
\(\xi^{D_+}\) and \(\xi^{D_-}\) are rational. It satisfies
\begin{equation}
 \frac{\xi(w+1)^{D_+}-\xi(w)^{D_+}}
      {\xi(w+1)^{D_-}-\xi(w)^{D_-}}
 \frac{\xi(w-1)^{D_+}-\xi(w)^{D_+}}
      {\xi(w-1)^{D_-}-\xi(w)^{D_-}}
 =\left(\frac{D_+}{D_-}\right)^2
      \xi(w)^{2(D_+-D_-)}.
 \label{shape:profile-reflection}
\end{equation}
The corresponding assertion for a pole holds after interchanging
\(Z\) and \(W\).
\end{lemma}

\begin{proof}
Write \(z=s-p\). By \eqref{shape:normal-reconstruction},
\(f=n/(\ii z)+O(1)\). Integration of
\eqref{shape:coordinate-derivatives} therefore gives
\begin{equation}
 Z(s)-Z(p)\sim c_Z z^{n+3},\qquad
 W(s)-W_0\sim c_W z^{3-n},
 \label{shape:center-powers}
\end{equation}
with nonzero constants \(c_Z,c_W\). Here \(W_0=W(p)\) when \(W\) is
finite, and any fixed constant may be used when \(W\) has a pole. If
\(n=3\), the leading term of \(W'\) is a nonzero multiple of \(z^{-1}\),
which cannot be the derivative of a meromorphic function. This proves
\(n\ne3\).

At sufficiently close punctured points of the zero-parameter slice,
\(\Gamma\) is finite, regular and nonisotropic. Its local inverse,
together with the rank-one identity for \(X\), gives
\begin{equation}
 X(y,t)=\Gamma(v(y,t)),\qquad
 v(y,t)=y+\sum_{j\geq1}t^{2j}v_j(y).
 \label{shape:local-impact-series}
\end{equation}
The even powers follow from evenness of \(X\). Possible joint poles or
indeterminacies at points where the zero-slice restrictions are regular
form a discrete set on that slice. After shrinking the punctured disk,
there are no such exceptions. The coefficients in
\eqref{shape:local-impact-series} glue there. Recursively inverting one
of the coordinates of \(\Gamma\) shows that every \(v_j\) is
meromorphic at \(p\): at each order the new coefficient is divided only
by a nonzero meromorphic derivative of that coordinate.

To take a leading limit with \(z=s-p\) and \(t\) on the same scale,
meromorphy of the coefficients alone is not enough. We need the pole bound
\begin{equation}
 \ord_p v_j\geq1-2j.
 \label{shape:impact-coefficient-order}
\end{equation}
Under \(z=tw\), this keeps every formal term \(t^{2j}v_j(p+tw)\)
at order at least \(t\), just as for the leading displacement \(z\).
We prove the bound using the arc-minus-chord action, divided by \(c\),
between shape parameters \(s\) and \(s+h\):
\begin{equation}
 \mathcal D(s,s+h)=
 \int_s^{s+h}f^{-2}
 -\sqrt{\left(\int_s^{s+h}f^{-2}a\right)
              \left(\int_s^{s+h}f^{-2}a^{-1}\right)}
 =\frac{h^3}{24}+O(h^5).
 \label{shape:point-action}
\end{equation}
The square root initially has the arc-length branch. To check the cubic
coefficient, expand each integral at its midpoint. The difference
between the normalized second derivative of \(f^{-2}\) and the average
of those of \(f^{-2}e^{\pm\ii\phi}\) is \(f^2\). Multiplication by
\(f^{-2}\) leaves the constant coefficient \(1/24\).

Write \(\mathcal D(s,s+h)=\sum_m d_m(s)h^m\). Substituting
\(s=p+z\), \(h=z\eta\), the two chord integrals factor as
\(z^{3+n}\eta\) and \(z^{3-n}\eta\), respectively, times functions
holomorphic near \((z,\eta)=(0,0)\) and nonzero there after division by
\(\eta\). Their product has the square root \(z^3\eta\) times a
holomorphic unit, with the chosen glancing sign. Thus
\begin{equation}
 \ord_p d_m\geq3-m.
 \label{shape:action-coefficient-order}
\end{equation}
This argument concerns short segments away from the center; the
resulting coefficient identities are meromorphic identities at the
center.

Reflection implies stationarity of the sum of two consecutive actions
\(\mathcal D\) on the lifted impacts in
\eqref{shape:local-impact-series}. The arc terms telescope, and the
chosen square roots have velocities tending to the same unit tangent,
so their signs give the ordinary reflected branch. Divide this
stationary equation by \(t^3\). For a general function \(v\), its
leading expression is
\[
 -\frac14v_yv_{yy}.
\]
Indeed the cubic action contributes one eighth of the difference of
the squared previous and next increments. At order \(t^{2j}\), the new
coefficient consequently enters as \(-v_j''/4\).

Assign weight one to each of \(z,t\). Assuming
\eqref{shape:impact-coefficient-order} at preceding orders, the formal
truncation of \(v-p\) through order \(t^{2j-2}\) and its increments
have lower total degree at least one.
Taylor shifts by \(\pm t\) preserve this property. Composition of
\(d_m\) with \(v\) preserves its lower degree \(3-m\), because each
derivative lowers the allowed order by at most one. The endpoint
derivatives of the unnormalized action therefore have total lower
degree at least two. After division by \(t^3\), the forcing at order
\(t^{2j}\) has order at least \(-1-2j\) in \(z\). Only finitely many
terms enter any fixed coefficient, since every step has positive
\(t\)-order. Hence \(\ord_p v_j''\geq-1-2j\), which proves
\eqref{shape:impact-coefficient-order} because \(v_j\) is meromorphic.

Let \(P(z,t)\) be a constant-coefficient polynomial combination of the
affine coordinates of \(X(p+z,t)\). Suppose that its zero-slice
restriction has order \(b\in\Z\) and leading coefficient \(c_P\ne0\).
Composition with \eqref{shape:local-impact-series} and
\eqref{shape:impact-coefficient-order} show that its coefficient of
\(t^l\) has order at least \(b-l\). Consequently
\[
 z^{-b}P(z,z\lambda)
\]
is holomorphic near \((z,\lambda)=(0,0)\). Here is the convergence
point in this assertion. The expression is a meromorphic germ and its
formal expansion in \(\lambda\), with Laurent coefficients in \(z\),
belongs to \(\C[[z,\lambda]]\) by the order bound. A meromorphic germ
whose quotient is formally a power series is holomorphic: apply
Weierstrass preparation and division, after a generic linear change of
coordinates, to its numerator and denominator. More explicitly, replace
the denominator by a Weierstrass polynomial times a unit, and divide
the analytic numerator by that polynomial. The analytic remainder has
degree below the polynomial degree. Uniqueness of formal division,
compared with the given formal quotient, makes its formal remainder
zero. The analytic remainder is consequently zero as well, and the
analytic quotient is the desired germ; see
\cite[Theorem~6.1.1]{Hormander1990}. Thus formal nonnegative orders
imply genuine holomorphy in this instance.

Its restriction to \(z=0\) is rational in \(\lambda\), since it is the
ratio of the first nonzero homogeneous parts of a local numerator and
denominator after the substitution \(t=z\lambda\). Formally that
restriction equals
\[
 c_P\mathcal U(\lambda)^b,\qquad \mathcal U(0)=1,
\]
where \(\mathcal U\) collects the total-degree-one terms of \(v-p\).
Taking \(P=Z-Z(p)\), for which \(b=D_+>0\), supplies an analytic
\(D_+\)-th root near one and proves convergence of \(\mathcal U\).
Set \(\xi(w)=w\mathcal U(1/w)\). For sufficiently large \(R_0\),
\begin{equation}
 t^{-b}P(tw,t)\longrightarrow c_P\xi(w)^b
 \label{shape:rational-leading-limit}
\end{equation}
locally uniformly, with derivatives, on \(|w|>R_0\); the limit is
rational. Applying this to the two coordinates in
\eqref{shape:center-powers} proves rationality of the stated powers.

In isotropic coordinates reflection says that the product of the two
incident chord slopes, expressed as \(Z\)-increment divided by
\(W\)-increment, is the square of the tangent slope. Substitute
\eqref{shape:rational-leading-limit} at \(w,w-1,w+1\). The constants
\(c_Z/c_W\) cancel, and the derivative of \(\xi\) cancels from the
tangent ratio. This gives \eqref{shape:profile-reflection}, first near
infinity and then as an identity of rational functions. Interchanging
the isotropic coordinates treats a pole of \(a\).
\end{proof}

\begin{lemma}[Every isotropic center is simple]
\label{shape:simple-centers}
All zeros and poles of \(a\) in \(D_H\) are simple. At each such center
\(p\), the curve has the local symmetry
\begin{equation}
 \Gamma(2p-s)=\Gamma(s),\qquad a(2p-s)=-a(s).
 \label{shape:center-reflection}
\end{equation}
\end{lemma}

\begin{proof}
It suffices to consider a zero of order \(n\). Use the preceding
lemma and put
\[
 d=\gcd(D_+,|D_-|),\quad R=\xi^d,\quad
 \alpha=D_+/d,\quad\beta=D_-/d,\quad
 \rho(w)=\frac{R(w+1)}{R(w)}.
\]
Bezout's identity makes \(R\) rational from the rationality of the two
powers of \(\xi\). Since \(R(w)\sim w^d\),
\(\rho(w)=1+d/w+O(w^{-2})\). Thus \(\rho\) is a nonconstant rational
map of some degree \(s\geq1\), with \(\rho(\infty)=1\). Define
\[
 F(x)=\frac{x^\alpha-1}{x^\beta-1}.
\]
After cancellation of the common powers of \(R(w)\),
\eqref{shape:profile-reflection} becomes
\begin{equation}
 F(\rho(w+1))=
 \frac{(\alpha/\beta)^2}{F(1/\rho(w))}.
 \label{shape:rational-shift-equation}
\end{equation}
We use the Riemann--Hurwitz formula: a rational map of degree
\(s\) has total ramification \(2s-2\); see \cite{Forster1981}.

If \(n=2\), then \(\alpha=5,\beta=1\), so \(F\) is a polynomial
of degree four with simple zeros at the four nontrivial fifth roots
of unity. At each preimage by \(\rho\) of one of these four values,
the right side of \eqref{shape:rational-shift-equation} has a pole
whose order is the local multiplicity of that preimage. On the left,
a pole has order divisible by four. Thus all these multiplicities
are divisible by four. Each of the four fibers contributes at least
\(3s/4\) to ramification, giving total at least \(3s>2s-2\), a
contradiction.

Suppose next that \(n>3\). Write \(\gamma=-\beta>0\) and
\(m=\alpha-\gamma=6/d\geq1\). The function
\[
 F(x)=-x^\gamma\frac{x^\alpha-1}{x^\gamma-1}
\]
has a pole of order \(\alpha\) at infinity. The function on the
right of \eqref{shape:rational-shift-equation}, considered as a
function of \(\rho\), has a pole of order \(\gamma\) at infinity
and simple poles at the nontrivial \(\alpha\)-th roots of unity.
These roots are distinct from its other possible exceptional values,
since \(\gcd(\alpha,\gamma)=1\).

Group the finitely many poles of \(\rho\) into maximal runs
\(w_0,w_0+1,\ldots,w_0+\ell\), with respective orders
\(p_0,\ldots,p_\ell\). Comparing pole orders at the points of the
run gives
\[
 \alpha p_{j+1}=\gamma p_j\quad(0\leq j<\ell).
\]
At the preceding point \(w_0-1\), which is not a pole of \(\rho\),
\(\rho\) must take a nontrivial \(\alpha\)-th root with multiplicity
\begin{equation}
 \alpha p_0=m\sum_{j=0}^\ell p_j+\gamma p_\ell.
 \label{shape:pole-run-seam}
\end{equation}
The preceding points of different runs are distinct. Their
ramification contributions are their displayed multiplicities minus
one. If \(m\geq2\), \eqref{shape:pole-run-seam} shows that these
points alone contribute at least \(ms\geq2s\), a contradiction.

If \(m=1\) and \(\gamma\geq2\), coprimality in the recurrence
implies \(\gamma^\ell\mid p_\ell\). Therefore
\(\gamma p_\ell-1\geq\ell+1\). If \(N_P\) is the number of
distinct poles, the preceding points contribute at least \(s+N_P\),
while the poles themselves contribute \(s-N_P\). Again the total is
at least \(2s\). The only remaining possibility is
\(m=\gamma=1\), \(\alpha=2\). All preceding points then take the
single value \(-1\), and their total multiplicity, by
\eqref{shape:pole-run-seam}, is strictly larger than \(s\). This is
impossible for one fiber of a degree-\(s\) map. Together with the
exclusion of \(n=3\), these arguments leave only \(n=1\).

For \(n=1\) we have \(d=2,\alpha=2,\beta=1\). Equation
\eqref{shape:rational-shift-equation} reads
\[
 (1+\rho(w+1))(1+1/\rho(w))=4.
\]
Equivalently, \(2/(\rho-1)\) increases by one under translation.
A rational one-periodic function is constant, so
\[
 \frac{2}{\rho(w)-1}=w+b_0,
 \qquad R(w)=(w+b_0)(w+b_0+1).
\]
The second identity follows by comparing the shift ratios and the
asymptotic \(R(w)\sim w^2\). In particular \(\xi\) is not rational,
since \(\xi^2\) has two simple zeros.

Choose a local coordinate \(\zeta=z+O(z^2)\) such that
\(W-W_0=c_W\zeta^2\). If the expansion of \(Z\) in \(\zeta\)
had a first odd term of order \(j\), subtracting the preceding even
terms as a polynomial in \(W-W_0\) would give a polynomial
combination \(P\) with order \(j\). By
\eqref{shape:rational-leading-limit}, \(\xi^j\) would be rational.
Since \(j\) is odd and \(\xi^2\) is rational, Bezout's identity would
make \(\xi\) rational, a contradiction. Thus both coordinates of
\(\Gamma\) are even in \(\zeta\). The involution
\(\zeta\mapsto-\zeta\) induces a local symmetry \(\sigma\) of the
shape parameter, fixing \(p\) with derivative \(-1\).

By \eqref{shape:intrinsic}, \(\sigma'^6=1\) on regular points.
Its derivative is therefore constant, and \(\sigma(s)=2p-s\).
The tangent-ratio identity gives \(a(\sigma(s))^2=a(s)^2\).
The simple zero fixes the sign to be negative. This proves
\eqref{shape:center-reflection}; the pole case is identical after
interchanging \(Z,W\).
\end{proof}

\begin{theorem}[Classification in the presence of isotropic centers]
\label{thm:isotropic-classification}
If \(a\) has a zero or a pole in \(D_H\), the boundary is an ellipse.
\end{theorem}

\begin{proof}
A center is nonreal, since \(|a|=1\) on the real axis. If \(H<\infty\),
the symmetry \eqref{shape:center-reflection} extends by identity on
the overlap of \(D_H\) with its reflection about the center. That
overlap is connected. The reflected strip extends past one of the
edges of \(D_H\), giving a meromorphic continuation of \(\Gamma\)
through an edge arc. This contradicts
\Cref{cor:shape-natural-boundary}. Hence \(H=\infty\).

Real symmetry supplies the conjugate center, and composition of the
two center reflections supplies a nonreal period, in addition to
\(2\pi\). Thus \(a\) is a nonconstant elliptic function. Let
\(\Lambda\) be its full period lattice. Composing reflections about
any two centers shows that twice their difference belongs to
\(\Lambda\). Consequently all zeros and poles lie in the four
classes of one translate of \(\frac12\Lambda/\Lambda\). They are
all simple by \Cref{shape:simple-centers}. The degree of \(a\) on
\(\C/\Lambda\) is therefore at most two. It cannot be one, because
a degree-one meromorphic map from a torus to the sphere would be a
biholomorphism. Its degree is exactly two.

Riemann--Hurwitz now gives four simple ramification points. Their
branch values \(\beta_1,\ldots,\beta_4\) are distinct and finite:
a fiber of degree two cannot contain two ramification points, and
the poles of \(a\) are simple. The quotient
\[
 \frac{a'^2}{\prod_{j=1}^4(a-\beta_j)}
\]
has neither zeros nor poles on the torus. At a ramification point
both numerator and denominator have order two; at a simple pole
both have order minus four. Hence the quotient is a nonzero
constant. We have obtained
\begin{equation}
 a'^2=P(a),\qquad \deg P=4.
 \label{shape:elliptic-normal-equation}
\end{equation}
Differentiating \(a(2p-s)=-a(s)\) gives
\(a'(2p-s)=a'(s)\), so \(P(-a)=P(a)\). Thus \(P\) is even.
On the real curve,
\[
 f^2=-\frac{a'^2}{a^2}
   =C+B\cos(2\phi)+D\sin(2\phi)>0
\]
for real constants \(B,C,D\). Therefore
\(f^2=n_\phi^{\mathsf T}Nn_\phi\) for a real positive-definite
symmetric matrix \(N\). Positivity follows on all directions since
\(\phi\) runs once around the circle.

For \(q_N(\phi)=n_\phi^{\mathsf T}Nn_\phi\), direct differentiation
gives
\[
 \left(\sqrt{q_N}\right)''+\sqrt{q_N}
 =\frac{\det N}{q_N^{3/2}}.
\]
For example, use \(q_N'=2e_\phi^{\mathsf T}Nn_\phi\) and
\(q_N''=2\operatorname{tr}N-4q_N\), and evaluate the determinant of
\(N\) in the orthonormal basis \((n_\phi,e_\phi)\). It follows that
the support
\[
 h_0(\phi)=\frac{c}{\det N}\sqrt{n_\phi^{\mathsf T}Nn_\phi}
\]
has curvature radius \(c/f^3\), the radius of the given boundary.
This is the support of an ellipse. The difference \(h-h_0\) solves
\((h-h_0)''+(h-h_0)=0\), so it is a linear combination of
\(\cos\phi,\sin\phi\), precisely a translation of the support.
\end{proof}

\subsection{A global lift when the normal is nonisotropic}

For the rest of this section suppose that \(a\) has neither zeros
nor poles in \(D_H\). Since the strip is simply connected, it has a
holomorphic logarithm. Choose \(\phi\) consistently with the real
normal. Then \(\phi,f=\phi'\) are holomorphic, and
\[
 \phi(s+2\pi)=\phi(s)+2\pi.
\]
The affine curve can still have poles, caused by zeros of \(f\).
The following theorem controls the impact coordinate even there.

\begin{theorem}[The global shape lift]
\label{thm:shape-lift}
Unless the boundary is a circle, there is a single-valued holomorphic
function \(\mathfrak s\) on the joint tube
\begin{equation}
 \mathcal T=\left\{(y,t):t\text{ is in the parameter disk},\quad
 -h_X(\bar t)<\Im y<h_X(t)\right\}
 \label{shape:joint-tube}
\end{equation}
such that
\begin{equation}
 X(y,t)=\Gamma(\mathfrak s(y,t)),\qquad
 |\Im\mathfrak s(y,t)|<H.
 \label{shape:global-composition}
\end{equation}
It satisfies
\begin{equation}
 \mathfrak s(y,0)=y,\qquad
 \mathfrak s(y+2\pi,t)=\mathfrak s(y,t)+2\pi,\qquad
 \mathfrak s(y,-t)=\mathfrak s(y,t).
 \label{shape:lift-symmetries}
\end{equation}
It also obeys real symmetry,
\(\mathfrak s(\bar y,\bar t)=\overline{\mathfrak s(y,t)}\).
After shrinking the parameter disk, for every real \(t\) the real
map \(y\mapsto\mathfrak s(y,t)\) descends to an orientation-preserving
analytic circle diffeomorphism. The inequality in
\eqref{shape:global-composition} is void when \(H=\infty\).
\end{theorem}

\begin{proof}
Near real glancing, the regular parametrization of the curve and the
rank-one identity give the asserted lift. Its differential satisfies
\begin{equation}
 c^2\mathfrak s_y^6
  =\frac{\det(X_y,X_{yy})^4}{(X_y\mathbin{\cdot}X_y)^3},
 \qquad
 \frac{\mathfrak s_t}{\mathfrak s_y}
  =\frac{X_t\mathbin{\cdot}X_y}{X_y\mathbin{\cdot}X_y}.
 \label{shape:lift-differential}
\end{equation}
Remove the proper analytic sets on which the meromorphic quantities
needed here have poles, vanish in a denominator, or make the chosen
sixth root zero. The complement containing the initial germ is
connected: removing complex analytic sets of positive codimension
does not disconnect a domain. Along any path in this complement the
sixth root, the closed differential it determines, and its local
primitive continue. Closedness continues from the initial germ.
On every compact regular path the differential is finite, and its
\(y\)-component is nonzero.

The composition identity persists for as long as the continued
primitive lies in \(D_H\). If \(H<\infty\), it cannot first reach an
edge of that strip. At such a first point the primitive is still
finite and locally submersive. Complete it to local product
coordinates \((\mathfrak s,z)\), and restrict the meromorphic map
\(X\) to a generic constant-\(z\) section. The identity on the inner
side would then continue the primary \(\Gamma\) meromorphically
through its edge, contradicting
\Cref{cor:shape-natural-boundary}. Thus the primitive remains in
the strip throughout continuation. For \(H=\infty\) there is no edge
to encounter.

After return along a loop, the continued primitive has the form
\[
 \mathfrak s\longmapsto\omega\mathfrak s+b,
 \qquad \omega^6=1.
\]
It consequently induces the symmetry
\begin{equation}
 \Gamma(\omega s+b)=\Gamma(s)
 \label{shape:possible-monodromy}
\end{equation}
on the overlap of the two relevant strips. That overlap is
connected, being the intersection of two strips. If \(H<\infty\),
a nonreal \(\omega\) makes the transformed strip cross an edge of
\(D_H\). If \(\omega=\pm1\), the same is true when \(\Im b\ne0\).
In either case \eqref{shape:possible-monodromy} would give an edge
patch. Thus only \(\omega=\pm1\), \(b\in\R\) remain.

For \(H=\infty\), a nonreal \(\omega\) gives a nonreal period by
conjugating the real translation \(s\mapsto s+2\pi\) with the affine
symmetry in \eqref{shape:possible-monodromy}. A nonreal translation
is already such a period. If \(\omega=-1\) and \(\Im b\ne0\),
compose that reflection with its conjugate, which is also a symmetry
by reality, to obtain the nonreal period \(b-\bar b\).
A period of \(\Gamma\) is a period of \(a\) by
\eqref{shape:normal-reconstruction}. Such a nonreal period, together
with \(2\pi\), would make the entire holomorphic function \(a\)
doubly periodic and hence constant. This contradicts its real
winding. Again only the real possibilities remain.

Real embeddedness rules out a real reflection in
\eqref{shape:possible-monodromy} and makes a real translation an
integer multiple of \(2\pi\). We conclude that
\[
 v=e^{\ii\mathfrak s}
\]
is single-valued on the regular complement. When \(H<\infty\), both
\(v\) and \(1/v\) are locally bounded, in fact between the bounds
given by \(e^{\pm H}\). They extend holomorphically across the
removed analytic sets by the removable-singularity theorem. The
bounds remain strict by the maximum principle, since \(v\) is
nonconstant on the zero slice. Hence the extension is nonvanishing
and still has \(|\log|v||<H\).

For completeness we address the unbounded-strip case separately;
bounded removability is not available there. The differential
\(\dd\mathfrak s\) is already single-valued, because all its
monodromy is translational. It extends meromorphically by
\eqref{shape:lift-differential}. Indeed, if a single-valued
holomorphic function off an analytic set has a meromorphic sixth
power, multiply it by a local denominator of that power; the
result is bounded, since its sixth power is holomorphic, and is
therefore removable. Apply this to \(\mathfrak s_y\), then use the
second identity in \eqref{shape:lift-differential} for the other
component.

Write \(\Gamma(s)=G(e^{\ii s})\), where \(G\) is meromorphic on
\(\C^*\). At a generic smooth point of a possible polar divisor,
restrict to a transverse disk, with the center deleted. The
single-valued function \(v\) on this punctured disk cannot have an
essential singularity. If it did, the image of every punctured
neighborhood would be dense; choosing sequences approaching any
regular value \(w\in\C^*\) would force
\(G(w)\) to equal the projective limit of the meromorphic map
\(X=G(v)\). The nonconstant map \(G\) would be constant.

Suppose instead that \(v\) has a zero of finite order \(k\). In a
local coordinate, \(v(z)=z^k u(z)\), \(u(0)\ne0\). Meromorphic
affine coordinates of \(X\) grow at most like \(|z|^{-N}\).
Every sufficiently small nonzero \(w\) has preimages satisfying
\(|z|\asymp|w|^{1/k}\). The identity \(G(w)=X(z)\) gives
polynomial growth of the affine coordinates of \(G\) at zero, hence
meromorphic extension there. A pole of \(v\) gives the analogous
extension at infinity. Reality,
\(G(1/\bar w)=\overline{G(w)}\) in Euclidean coordinates, supplies
extension at the other end. Consequently \(G\) is rational.
Equations \eqref{shape:normal-reconstruction} then make \(a\)
rational in \(w=e^{\ii s}\). A rational function with no zeros or
poles on \(\C^*\) is \(Cw^k\); its real winding is one, so \(k=1\).
This makes \(f\) constant and the boundary a circle, the excluded
case.

It follows that \(v\) extends with a nonzero finite value on every
generic transverse disk. Therefore
\(\dd\mathfrak s=\dd v/(\ii v)\) has no pole on such a disk.
Generic transverse directions detect all leading polar coefficients
of a meromorphic one-form: a nonzero coefficient covector cannot
annihilate every direction having nonzero normal component.
Thus \(\dd\mathfrak s\) has no polar divisor. Removability in
codimension two makes it holomorphic throughout \(\mathcal T\).

Finally, \(\mathcal T\) is simply connected. Its vertical fibers in
\(y\) are intervals containing \(\Im y=0\), so vertical contraction
retracts it to \(\R\) times the parameter disk. For finite \(H\),
take a holomorphic logarithm of the nonvanishing extension of \(v\).
For infinite \(H\), take a primitive of the holomorphic closed
one-form \(\dd\mathfrak s\). Normalize by the original germ. This
gives the required single-valued lift. The composition and all
symmetries in \eqref{shape:lift-symmetries}, including real symmetry,
follow by identity from the initial domain. On the real phase circle,
\(\mathfrak s_y=1+O(t^2)\) uniformly. Shrinking the disk makes this
derivative positive for real \(t\); the degree-one periodicity then
gives the asserted circle diffeomorphism.
\end{proof}

\subsection{Zeros of the speed and the first displacement coefficient}

\begin{lemma}[Simple zeros and a holomorphic potential]
\label{lem:f-zeros}
Every zero of \(f\) in \(D_H\) is simple. At each such zero,
\(f''=0\). Consequently
\begin{equation}
 q=\frac{f''}{f}
 \label{shape:potential-definition}
\end{equation}
extends holomorphically throughout \(D_H\).
\end{lemma}

\begin{proof}
Suppose that \(f\) vanishes at \(s_0\) to order \(m\geq1\). Use
fixed tangent and normal coordinates for the angle \(\phi(s_0)\),
and write \(z=s-s_0\). Since
\(\phi(s)-\phi(s_0)\) has order \(m+1\), the tangent and normal
components of \(\Gamma'\) in
\eqref{shape:coordinate-derivatives} have leading orders
\(-2m\) and \(1-m\), respectively, with nonzero coefficients.
If \(m=2\), the normal derivative has a nonzero residue, contrary
to meromorphy. If \(m\geq3\), integration gives tangent and normal
coordinate poles of orders \(2m-1\) and \(m-2\).

By \Cref{thm:shape-lift}, the lift is holomorphic near \((s_0,0)\).
The equation \(\mathfrak s(y,t)=s_0\) defines
\(y=s_0+O(t^2)\). For sufficiently small \(t\ne0\), the two
neighboring shape coordinates at phase shifts \(\pm t\) are
\(s_0\pm t+O(t^2)\); their impacts are finite. Along the central
curve, the chord to the previous impact therefore has normal-to-tangent
slope of order at least \(m+1\). Its normal angle \(\theta\)
can be chosen holomorphically, and
\[
 d=\theta-\phi(\mathfrak s)
   =O\big((\mathfrak s-s_0)^{m+1}\big).
\]
For its signed offset \(p_{\rm chord}\), differentiation of
\(p_{\rm chord}=n_\theta\cdot\Gamma(\mathfrak s)\) gives
\begin{equation}
 \dd p_{\rm chord}\wedge\dd\theta
 =c f(\mathfrak s)^{-2}\sin d\,
       \dd\mathfrak s\wedge\dd d
 =\pm A(t)\,\dd y\wedge\dd t.
 \label{shape:chord-area}
\end{equation}
The final identity is the continued area identity
\eqref{eq:area-coefficient}; choosing a local orientation can change
only its sign. Since \(A(t)=ct/4+O(t^3)\), it is nonzero for small
\(t\ne0\).

Set \(z=\mathfrak s-s_0\) and write \(d=z^{m+1}b\).
Then \(\dd\mathfrak s\wedge\dd d
=z^{m+1}\dd\mathfrak s\wedge\dd b\). The left side of
\eqref{shape:chord-area} thus vanishes to order at least
\(-2m+(m+1)+(m+1)=2\), a contradiction. Hence \(m=1\).
At a simple zero write \(f=Az+Bz^2+O(z^3)\), \(A\ne0\).
Since \(a'=\ii fa\), we have \(a=a(s_0)+O(z^2)\). The residue
of \(a/f^2\) is \(-2a(s_0)B/A^3\). It vanishes because
\(\ii c a/f^2=Z'\) is a derivative of a meromorphic function.
Thus \(B=0\), which is \(f''(s_0)=0\). The asserted removability
of \(q\) follows.
\end{proof}

We now fix the harmless freedom to translate phase as a function of
\(t\). Let
\[
 b(t)=\frac{1}{2\pi}\int_0^{2\pi}
             (\mathfrak s(y,t)-y)\,\dd y.
\]
This is analytic and even, real for real \(t\), and vanishes at zero.
Precompose \(L,X,\mathfrak s\) with \(y\mapsto y-b(t)\), shrinking
the parameter disk if necessary. Periodicity and contour translation
then give
\begin{equation}
 \left\langle\mathfrak s(y,t)-y\right\rangle=0.
 \label{shape:mean-normalization}
\end{equation}
All rotation, area and symmetry identities are preserved. The two
widths acquire the same harmonic summand \(\Im b(t)\); hence the
width relation, superharmonic comparison and no-patch properties are
also preserved. We retain the previous notation for the normalized
data. Averaging on any horizontal row in the tube now gives
\(\langle\Im\mathfrak s\rangle=\Im y\). The shape constraint
therefore implies
\begin{equation}
 h_X(t)\leq H.
 \label{shape:width-bound}
\end{equation}

\begin{proposition}[The first impact displacement]
\label{shape:first-displacement}
For the normalization \eqref{shape:mean-normalization},
\[
 \mathfrak s(y,t)=y+t^2V(y)+O(t^4),
\]
where \(V\) is holomorphic and periodic on \(D_H\), has zero mean,
and satisfies
\begin{equation}
 V'=\frac{f^2+8q-\langle f^2+8q\rangle}{120}.
 \label{eq:V-coefficient}
\end{equation}
\end{proposition}

\begin{proof}
It suffices to calculate on real nonsingular data and then continue
by identity. In \eqref{shape:point-action} put
\(g=f^{-2}\), \(g_\pm=g e^{\pm\ii\phi}\). Midpoint integration
through the fourth derivative gives
\begin{equation}
 \mathcal D(s,s+h)
 =\frac{h^3}{24}+d_5(s+h/2)h^5+O(h^7).
 \label{shape:fifth-action}
\end{equation}
For clarity, define
\[
 B_j=\frac12\left(\frac{g_+^{(j)}}{g_+}
                         +\frac{g_-^{(j)}}{g_-}\right),\qquad
 C_2=2(g'/g)f+f'=-3f'.
\]
Multiplying the two midpoint integral expansions and taking the
glancing square root gives
\[
 d_5=g\left[
       \frac{g''''/g-B_4}{1920}
       -\frac{C_2^2}{2\cdot24^2}\right].
\]
Differentiating \(g_\pm\) explicitly yields
\[
 g''''/g-B_4=-8ff''+15(f')^2-f^4.
\]
The two \((f')^2\) contributions cancel, since
\(15/1920=9/(2\cdot24^2)\). Therefore
\begin{equation}
 d_5=-\frac{f^2+8q}{1920}.
 \label{shape:d5}
\end{equation}

At the middle impact \(s=\mathfrak s(y,t)\), the previous and next
shape increments are
\[
 h_-=t+t^3V'-\tfrac12t^4V''+O(t^5),\qquad
 h_+=t+t^3V'+\tfrac12t^4V''+O(t^5).
\]
The cubic term in the stationary action contributes
\(-t^5V''/4\). The fifth-order terms contribute
\(-5t^5d_5'+t^5d_5'\): the first part comes from differentiating
the increments and comparing the two midpoint values, and the
second from differentiating the midpoints. Hence stationarity at
order \(t^5\) reads
\[
 -V''/4-4d_5'=0.
\]
Substitute \eqref{shape:d5}, integrate once, and use periodicity of
\(V\) to determine the integration constant. This proves
\eqref{eq:V-coefficient}. The zero mean of \(V\) follows from
\eqref{shape:mean-normalization}.
\end{proof}
